%!TEX TS-program = xelatex
%!TEX encoding = UTF-8 Unicode
%-------------------------------------------------------------------------------
%  ĐƠN XIN VIỆC (COVER LETTER) - VÕ VƯƠNG KHÁNH VY
%
%  Mẫu: Awesome-CV của Claud D. Park (posquit0) - CC BY-SA 4.0
%       https://github.com/posquit0/Awesome-CV
%
%  Biên dịch:  make        (hoặc: xelatex don-xin-viec.tex, chạy 2 lần)
%
%  >>> QUAN TRỌNG: đơn này đang viết cho HỆ THỐNG NHÀ THUỐC TRUNG SƠN.
%      Mỗi lần nộp cho công ty khác PHẢI sửa lại 4 chỗ được đánh dấu
%      [SỬA CHO TỪNG CÔNG TY] bên dưới. Gửi đơn ghi sai tên công ty là
%      bị loại ngay. Xem thêm HUONG-DAN-UNG-TUYEN.md.
%-------------------------------------------------------------------------------

\documentclass[11pt, a4paper]{awesome-cv}

% Thiết lập chung + thông tin liên hệ (dùng chung với main.tex)
%-------------------------------------------------------------------------------
%  CẤU HÌNH DÙNG CHUNG cho CV (main.tex) và Đơn xin việc (don-xin-viec.tex)
%
%  Sửa thông tin liên hệ ở đây MỘT LẦN là cả hai tài liệu cùng cập nhật.
%-------------------------------------------------------------------------------

\geometry{left=1.4cm, top=.8cm, right=1.4cm, bottom=1.4cm, footskip=.5cm}

% Màu chủ đạo. Các lựa chọn có sẵn của Awesome-CV: awesome-emerald, awesome-skyblue,
% awesome-red, awesome-pink, awesome-orange, awesome-nephritis, awesome-concrete,
% awesome-darknight. Ở đây dùng màu xanh y tế cho phù hợp ngành Dược.
\definecolor{awesome}{HTML}{00796B}

%-------------------------------------------------------------------------------
% SỬA MẪU GỐC CHO HỢP VỚI TIẾNG VIỆT
%-------------------------------------------------------------------------------
% Mẫu gốc tô màu 3 KÝ TỰ đầu của tiêu đề mục (\StrSplit{...}{3}). Với tiếng Việt,
% cách này cắt ngang giữa từ ("Kin|h nghiệm", "Chứ|ng chỉ") trông như lỗi hiển thị,
% nên ở đây tô màu toàn bộ tiêu đề — cho cả CV lẫn đơn xin việc.
\setbool{acvSectionColorHighlight}{false}
\renewcommand*{\sectionstyle}[1]{{\color{awesome}\sectionstyleface{#1}}}
\renewcommand*{\lettersectionstyle}[1]{{\color{awesome}\lettersectionstyleface{#1}}}

% Mẫu gốc dùng chữ small-caps (\scshape) cho dòng chức danh, chân trang, tên vị trí
% và địa chỉ người nhận. Font không có glyph small-caps cho các chữ tiếng Việt có dấu
% (ư, ĩ, ẳ, ơ...) nên chúng rơi về chữ thường, làm dòng chữ cao thấp không đều.
% Thay bằng chữ IN HOA để giữ đúng ý đồ thiết kế mà vẫn hiển thị đủ dấu.
\renewcommand*{\headerpositionstyle}[1]{{\fontsize{7.6pt}{1em}\bodyfont\color{awesome}\MakeUppercase{#1}}}
\renewcommand*{\footerstyle}[1]{{\fontsize{8pt}{1em}\footerfont\color{lighttext}\MakeUppercase{#1}}}
\renewcommand*{\entrypositionstyle}[1]{{\fontsize{8pt}{1em}\bodyfont\color{graytext}\MakeUppercase{#1}}}
\renewcommand*{\subentrypositionstyle}[1]{{\fontsize{7pt}{1em}\bodyfont\color{graytext}\MakeUppercase{#1}}}
% Địa chỉ người nhận giữ nguyên chữ thường cho dễ đọc (tên công ty, đường, phường)
\renewcommand*{\recipientaddressstyle}[1]{{\fontsize{9pt}{1em}\bodyfont\color{graytext} #1}}

% Thu gọn khoảng cách giữa các mục để tài liệu vừa đúng 1 trang
\renewcommand{\acvSectionTopSkip}{2.4mm}
\renewcommand{\acvSectionContentTopSkip}{2mm}

\renewcommand{\acvHeaderSocialSep}{\quad\textbar\quad}

% Ngày tháng theo định dạng Việt Nam (\today của LaTeX trả về tiếng Anh)
\makeatletter
\newcommand{\ngayhomnay}{\two@digits{\day}/\two@digits{\month}/\the\year}
\makeatother

%-------------------------------------------------------------------------------
% THÔNG TIN CÁ NHÂN
%-------------------------------------------------------------------------------
\photo[circle,edge,left]{./avatar}
\name{Võ Vương}{Khánh Vy}
\position{Dược sĩ Cao đẳng~~·~~Dược sĩ tư vấn}
\address{KV. Phụng Thạnh 1, P. Thuận Hưng, TP. Cần Thơ}

\mobile{0949 507 278}
\email{vyk5045@gmail.com}
\dateofbirth{20/10/2005}

\quote{“Tư vấn đúng thuốc, đúng liều, đúng người --- để mỗi khách hàng quay lại vì tin tưởng.”}


% Bỏ câu trích dẫn ở đầu trang: CV đã có, còn lá đơn thì chính nội dung thư đã nói
% thay. Bỏ đi cũng giúp toàn bộ lá đơn vừa đúng một trang.
\makeatletter
\let\@quote\undefined
\makeatother

\hypersetup{
  pdftitle    = {Đơn xin việc - Võ Vương Khánh Vy - Dược sĩ tư vấn},
  pdfauthor   = {Võ Vương Khánh Vy},
  pdfsubject  = {Đơn xin việc vị trí Dược sĩ tư vấn tại TP. Cần Thơ},
  pdfkeywords = {đơn xin việc, dược sĩ tư vấn, nhà thuốc, GPP, Cần Thơ},
  pdflang     = {vi-VN},
}

%-------------------------------------------------------------------------------
% THÔNG TIN LÁ ĐƠN
%-------------------------------------------------------------------------------
% [SỬA CHO TỪNG CÔNG TY] (1/4) — tên bộ phận và địa chỉ nơi nhận
\recipient
  {Bộ phận Tuyển dụng}
  {Hệ thống Nhà thuốc Trung Sơn\\Công ty TNHH Trung Sơn Alpha\\TP. Cần Thơ}

\letterdate{Cần Thơ, ngày \ngayhomnay}

% [SỬA CHO TỪNG CÔNG TY] (2/4) — tên vị trí ứng tuyển
\lettertitle{Đơn xin việc --- Vị trí Dược sĩ tư vấn}

\letteropening{Kính gửi Quý Công ty,}
\letterclosing{Trân trọng,}
\letterenclosure[Hồ sơ đính kèm]{CV chi tiết, bằng tốt nghiệp, bảng điểm, giấy chứng nhận thời gian thực hành, CCCD}

% Mẫu gốc chừa 3 dòng trống giữa tên người ký và dòng hồ sơ đính kèm. Thu lại còn
% 2 dòng (vẫn đủ chỗ ký tay) để cả lá đơn vừa đúng một trang.
\makeatletter
\renewcommand*{\makeletterclosing}{%
  \vspace{2mm}
  \lettertextstyle{\@letterclosing} \\\\
  \letternamestyle{\@firstname\ \@lastname}
  \ifthenelse{\isundefined{\@letterenclosure}}
    {\\}
    {%
      \\\\
      \letterenclosurestyle{\@letterenclname: \@letterenclosure}%
    }
}
\makeatother

%-------------------------------------------------------------------------------
\begin{document}

\makecvheader[R]

\makecvfooter
  {Cần Thơ, \ngayhomnay}
  {Võ Vương Khánh Vy~~~·~~~Đơn xin việc}
  {}

\makelettertitle

%-------------------------------------------------------------------------------
% NỘI DUNG LÁ ĐƠN
%-------------------------------------------------------------------------------
\begin{cvletter}

\lettersection{Về tôi}
Tôi là Võ Vương Khánh Vy, Dược sĩ Cao đẳng vừa tốt nghiệp loại Giỏi hệ chính quy tại
Trường Cao đẳng Y tế An Giang. Gia đình tôi có một nhà thuốc đạt chuẩn GPP tại
P. Thốt Nốt, TP. Cần Thơ, nên trong suốt quá trình học tôi vừa học vừa phụ việc và thực
hành chuyên môn tại đây --- nhờ vậy, những gì học trên giảng đường đều được nhìn lại
hằng ngày qua công việc thực tế tại quầy thuốc.

Công việc tôi phụ trách là hướng dẫn khách hàng sử dụng thuốc OTC dưới sự hướng dẫn
trực tiếp của dược sĩ phụ trách chuyên môn, làm quen với danh mục thuốc và thực phẩm
chức năng thông dụng, sắp xếp và bảo quản thuốc đúng điều kiện, kiểm tra hạn dùng trên
quầy theo quy định GPP, cùng việc tra cứu thuốc và lập hoá đơn trên phần mềm quản lý
nhà thuốc.

% [SỬA CHO TỪNG CÔNG TY] (3/4) — thay tên công ty và lý do chọn nơi này
\lettersection{Vì sao tôi chọn Nhà thuốc Trung Sơn}
Làm ở nhà thuốc gia đình cho tôi sự quen thuộc với nghề, nhưng tôi hiểu rằng để tiến xa
hơn mình cần một môi trường có quy trình chuẩn hoá và được đào tạo bài bản. Nhà thuốc
Trung Sơn là hệ thống uy tín lâu năm tại khu vực Đồng bằng sông Cửu Long, đề cao vai trò
tư vấn chuyên môn của dược sĩ --- đây là nơi tôi tin mình sẽ học được nhiều nhất trong
những năm đầu đi làm. Tôi cũng sinh sống tại TP. Cần Thơ nên thuận tiện đi lại giữa các
chi nhánh và sẵn sàng nhận ca cuối tuần, ngày lễ.

\lettersection{Vì sao là tôi}
Tôi không đến với nghề từ con số không: ba năm phụ việc tại quầy thuốc đã cho tôi sự
quen thuộc với môi trường nhà thuốc, với cách trò chuyện và lắng nghe khách hàng, cũng
như với nhịp độ của những lúc đông khách. Tôi tin mình sẽ tiếp cận công việc nhanh và
cần ít thời gian làm quen hơn một bạn mới hoàn toàn.

Điều tôi mang theo từ những năm đó là ý thức rằng nghề dược đặt sự an toàn của người
bệnh lên trước hết: hỏi kỹ trước khi đưa thuốc, không tư vấn ngoài phạm vi hiểu biết
của mình, và luôn hỏi lại dược sĩ phụ trách khi chưa chắc chắn. Tôi cũng là người cẩn
thận, chịu khó và sẵn sàng học từ những việc nhỏ nhất.

Tôi đang tích lũy đủ thời gian thực hành để đề nghị cấp Chứng chỉ hành nghề dược và dự
định học liên thông Đại học Dược trong 2--3 năm tới; tôi tìm một nơi để gắn bó lâu dài,
không phải một chỗ làm tạm thời. Rất mong có cơ hội được trao đổi trực tiếp với Quý Công
ty --- xin chân thành cảm ơn.

\end{cvletter}

%-------------------------------------------------------------------------------
\makeletterclosing

\end{document}
